\documentclass[conference]{IEEEtran}

% ── Packages ────────────────────────────────────────────────────────────────
\usepackage[utf8]{inputenc}
\usepackage[T1]{fontenc}
\usepackage[english]{babel}
\usepackage{graphicx}
\usepackage{hyperref}
\usepackage{booktabs}
\usepackage{array}
\usepackage{amsmath}
\usepackage{amssymb}
\usepackage{cite}
\usepackage{enumitem}
\usepackage{xcolor}

\hypersetup{colorlinks=true, linkcolor=blue, urlcolor=blue, citecolor=blue}

% ════════════════════════════════════════════════════════════════════════════
\begin{document}

\title{Sprint 4 Results Evaluation:\\
Puzzlebot Challenge --- TE3003B}

\author{
\IEEEauthorblockN{
Juan José Jáuregui Barba A00836722 \\
Victor Alejandro Meneses Garza A01384002 \\
Rosendo De Los Rios Moreno A01198515 \\
Jordan Palafox Salinas A00835705 \\
Hugo Daniel Castillo Ovando A00836025
}
}
\maketitle

% ── Abstract ─────────────────────────────────────────────────────────────────
\begin{abstract}
This report documents the progress of Sprint~4 of the TE3003B
\textit{Puzzlebot Challenge} project. Building on the components delivered
in Sprint~3 (EKF, occupancy grid, SLAM prototype, A* + Bug1 navigation in
simulation, YASMIN state machine skeleton, Flask dashboard prototype, and
LPC + VQ voice module), Sprint~4 closed the remaining gaps toward real-robot
deployment: the SLAM and real-world mapping stack was debugged and validated
on hardware (DDS, clock offsets, QoS, TF pipeline); a rolling LiDAR local
costmap with temporal decay was added to navigation; the dashboard received
live map rendering and voice-command dispatch; the HMM speech recognizer path
was fixed and ten per-word models trained; the Forklift FSM was packaged as
an independent submodule; the lifter mechanism was redesigned with a
servo + FPGA assembly and validated end-to-end on the physical robot; and the
URDF model, lifter meshes, and Gazebo simulation world were finalized.
The complete architecture now runs end-to-end in simulation and the
real-world mapping and lifting pipelines have been validated on the Jetson Nano.
\end{abstract}

\begin{IEEEkeywords}
ROS~2, Puzzlebot, EKF, ICP, SLAM, A*, Bug1, local costmap, QR alignment,
pure pursuit, Bézier, YASMIN, FPGA, ArUco, HMM, speech recognition,
dashboard, Flask, SocketIO
\end{IEEEkeywords}

% ════════════════════════════════════════════════════════════════════════════
\section{Introduction}

Sprint~4 was the hardware-integration and closing sprint. Sprint~3 delivered
functional components validated in simulation: an EKF with ICP corrections,
an occupancy grid built from simulated scans, A* + Bug1 + PID navigation
tested in Gazebo, a YASMIN state machine with the full state topology, an
HMM-based voice node, a Flask dashboard prototype, and the forklift HAL with
Jetson.GPIO. The goal of Sprint~4 was to bring each of those components to
hardware readiness and to close the remaining open items.

The team's core design decision remained unchanged:
\textbf{the entire stack is a custom implementation}. Nav2, Whisper, and
\texttt{robot\_localization} are not used.


\begin{table}[h]
\caption{Tasks planned on the Sprint 4 board}
\label{tab:sprint4_board}
\renewcommand{\arraystretch}{1.2}
\centering
\begin{tabular}{>{\raggedright\arraybackslash}p{0.85\columnwidth}}
\toprule
\textbf{Task} \\
\midrule
Full system integration test on real hardware \\
QR-based visual alignment (docking to pallet position) \\
Real-world mapping stack on Jetson Nano \\
Local costmap with temporal decay \\
Dashboard: live occupancy-grid rendering \\
Dashboard: voice-command dispatch \\
HMM voice recognizer fix + 11-word model training \\
Forklift FSM submodule \\
SPI Jetson--FPGA validation \\
ArUco--EKF bridge node \\
CNN trailer detection \\
\bottomrule
\end{tabular}
\end{table}

% ════════════════════════════════════════════════════════════════════════════
\section{Sprint 4 Progress by Deliverable}

\subsection{M1. Industrial Standards in Robotics}

Documentation was updated to cover two additional standards relevant to
mobile robots that operate alongside humans:

\paragraph{ISO 10218-1/2 (Industrial robot safety)} the standard defines
safety requirements for industrial robot systems and integration. Applied to
the Puzzlebot:
\begin{itemize}[noitemsep, leftmargin=*]
  \item \textbf{Speed and separation monitoring:} maximum translational
        velocity capped at 0.3~m/s via \texttt{nav\_params.yaml}; the
        Bug1 node enforces a minimum separation distance from obstacles.
  \item \textbf{Protective stop:} the \texttt{stop} voice token and the
        dashboard emergency button publish zero velocity to \texttt{/cmd\_vel}
        and raise \texttt{stop\_flag} in the YASMIN blackboard within one
        control cycle.
\end{itemize}

\paragraph{ISO 13482 (Personal care robots — safety)} although the Puzzlebot
is a warehouse AMR, ISO~13482 provides the risk-assessment framework for
mobile service robots. The OWL ontology completed in Sprint~3 was extended
with classes \texttt{HazardZone} and \texttt{SafetyConstraint} linked to the
robot's workspace zones.

\subsection{M2. FPGA for the Forklift }

The \texttt{lifting} package from Sprint~3 was validated on hardware and the
Forklift FSM was packaged as an independent git submodule
(\texttt{src/forklift\_fsm}). Sprint~4 closes the loop: the full
servo + FPGA lifter assembly was mechanically redesigned, fabricated, and
demonstrated working end-to-end on the physical robot.

\paragraph{HAL validation} the \texttt{lifting\_node} was tested on the
Jetson Nano with the real FPGA. The three GPIO pins (BCM~11, 13, 15) are
written correctly by the ROS~2 node; the FPGA decodes the binary value and
drives the servo to the corresponding height level. Named level constants
in \texttt{lifting\_params.yaml}:

\begin{itemize}[noitemsep, leftmargin=*]
  \item \texttt{level\_rest = 0} — fully lowered
  \item \texttt{level\_transport = 1} — travel height
  \item \texttt{level\_pick\_floor = 3} — insert forks at floor level
  \item \texttt{level\_carry = 4} — carry height (after pick)
  \item \texttt{level\_pick\_rack2 = 5} — rack level~2 pick
  \item \texttt{level\_max = 7} — maximum height
\end{itemize}

\paragraph{Forklift FSM submodule} the Forklift finite-state machine was
separated into an independent repository and added as a git submodule at
\texttt{src/forklift\_fsm}. This allows the FPGA firmware team to develop
independently while the ROS~2 stack consumes the interface through the
\texttt{lifting\_node}.

\paragraph{New lifter mechanical design} the lifter mechanism was redesigned
from scratch for Sprint~4. The new design uses a servo motor driven by the
FPGA output and a custom-machined carriage to achieve the eight height levels
defined in Table~\ref{tab:lifter_encoding}. Figure~\ref{fig:lifter_cad} shows
the CAD model; Figure~\ref{fig:lifter_demo} links to the validation video.

% ── IMAGE 1 ─────────────────────────────────────────────────────────────────
% Place: img/lifter_cad.png
% Content: CAD render or screenshot of the new lifter design showing the
%          servo mount, carriage rail, and fork assembly.  Isometric view
%          preferred.  Export from your CAD tool at ≥1920 px wide.
\begin{figure}[htbp]
  \centering
  \includegraphics[width=\columnwidth]{img/lifter_cad.png}
  \caption{CAD model of the redesigned lifter assembly. The servo (driven by
           the FPGA output) translates the 3-bit GPIO command into one of
           eight discrete fork heights via a captive carriage mechanism.}
  \label{fig:lifter_cad}
\end{figure}

% ── VIDEO LINK ───────────────────────────────────────────────────────────────
% Replace the placeholder URL below with the actual Google Drive share link
% once the video is uploaded.
\begin{figure}[htbp]
  \centering
  \fbox{\parbox{0.9\columnwidth}{\centering\vspace{4pt}
    \textbf{Lifter demo video} \\[4pt]
    Servo + FPGA end-to-end test: ROS~2 publishes \texttt{/lifter\_level}
    commands 0--7; the Jetson writes the GPIO bits; the FPGA drives the servo
    through all eight height levels. \\[6pt]
    \href{https://drive.google.com/file/d/1Sq8gjUPFCLl14B59_lK673tai0RfwYfN/view?usp=sharing}{%
      \texttt{drive.google.com/lifter_servo_fpga}} \\[4pt]
  }}
  \caption{Lifter validation demo (Google Drive). The video shows the complete
           command chain: \texttt{lifting\_node} → GPIO → FPGA → servo →
           fork height.}
  \label{fig:lifter_demo}
\end{figure}

\begin{table}[h]
\caption{Binary encoding of the forklift (3 GPIO → FPGA)}
\label{tab:lifter_encoding}
\renewcommand{\arraystretch}{1.2}
\centering
\begin{tabular}{ccc>{\raggedright\arraybackslash}p{3.5cm}}
\toprule
\textbf{Level} & \textbf{GPIO} & \textbf{Use} \\
\midrule
0 & 000 & Rest / fully lowered \\
1 & 001 & Transport height \\
2 & 010 & Level 2 \\
3 & 011 & Pick pallet on floor \\
4 & 100 & Carry height \\
5 & 101 & Pick pallet rack level 2 \\
6 & 110 & Level 6 \\
7 & 111 & Maximum height \\
\bottomrule
\end{tabular}
\end{table}

\subsection{M3. Automatic Map Creation}

The real-world mapping stack was the first major hardware challenge of the
sprint. Three separate issues prevented the occupancy grid from being built
correctly on the Jetson:

\paragraph{DDS and clock offset} FastDDS required \texttt{IGN\_IP=localhost}
to avoid multicast discovery collisions with the simulation host. A persistent
clock offset between the Jetson system clock and the ROS~2 Time source caused
the TF tree to reject laser transforms; this was resolved by switching the
SLAM node to wall-clock time and adding a \texttt{transform\_tolerance} of
0.3~s to all TF lookups.

\paragraph{QoS mismatch} the LiDAR driver published \texttt{BEST\_EFFORT}
scans while the SLAM node expected \texttt{RELIABLE}; the resulting silent
drop of every scan meant the map never updated. The real-world launch file
was updated to match QoS profiles.

\paragraph{ICP improvements} the \texttt{icp.py} implementation received
major improvements: dynamic keyframe management, outlier
rejection with a distance threshold, and a minimum-correspondences guard that
prevents degenerate solutions in low-feature corridors.

\paragraph{Occupancy grid} \texttt{slam\_node.py}:
raycasting was rewritten with Bresenham's algorithm for speed, incremental
log-odds update was replaced with a probabilistic model, and a saved-map
loader was added so the robot can start a session with a pre-built map.
After these fixes, the Jetson Nano builds a clean, smear-free occupancy grid
from the RPLidar A1 in real-world conditions.

% ── IMAGE 2 ─────────────────────────────────────────────────────────────────
% Place: img/real_world_map.png
% Content: Screenshot of RViz2 showing the occupancy grid built on the
%          real Jetson Nano.  The warehouse perimeter walls and rack
%          obstacles should be clearly visible.  Side-by-side with the
%          Gazebo simulation map is ideal.
\begin{figure}[htbp]
  \centering
  \includegraphics[width=\columnwidth]{img/real_world_map.jpeg}
  \caption{Occupancy grid built on the Jetson Nano with the corrected
           real-world SLAM stack. Left: Gazebo simulation reference map.
           Right: map built from the physical RPLidar A1 scan in the lab.}
  \label{fig:real_map}
\end{figure}

\subsection{MR. Digital Twin --- Gazebo Simulation}

The warehouse Gazebo world was updated to match the physical competition
environment: rack models (\texttt{rack.dae}) were authored from scratch,
collision geometry was tuned in \texttt{model.sdf}, and the URDF xacro
was consolidated into a single \texttt{gazebo.launch.py} bringup file.
The digital twin allows full stack validation (SLAM, navigation, QR
alignment, lifter) without hardware access.

% ── IMAGE: digital twin ──────────────────────────────────────────────────────
% Place: img/gazebo_digital_twin.png
% Content: Gazebo GUI screenshot of the warehouse simulation world showing:
%   - The robot inside the warehouse floor plan
%   - At least one rack model with its collision mesh visible (View > Collisions)
%   - Ideally RViz2 side-by-side or inset: the occupancy grid the SLAM node
%     builds from simulated LiDAR in the same session
%   Capture at a top-down or 45° isometric camera angle so the layout is clear.
\begin{figure}[htbp]
  \centering
  \includegraphics[width=\columnwidth]{img/gazebo_digital_twin.jpeg}
  \caption{Gazebo digital twin of the competition warehouse. The scene
           includes the rack obstacles (\texttt{rack.dae}) and the
           Puzzlebot model. The inset shows the occupancy grid produced
           by the SLAM node from simulated LiDAR in the same session.}
  \label{fig:gazebo_twin}
\end{figure}

\subsection{M4. Voice Commands to Control the Robot}

The HMM-based speech recognizer inherited from Sprint~3 had a path
resolution bug that caused it to search for model files in a hardcoded
directory that does not exist after installation via \texttt{pip install -e}.
The fix resolves the model path using \texttt{importlib.resources} (or a
\texttt{pkg\_resources} fallback), making the node portable across any
workspace layout.

Eleven HMM models (\texttt{.npz}) were trained and committed:
\texttt{start, stop, forward, backward, left, right, lift, leave,
break, continue, resume}.

\paragraph{Recognition pipeline}
\[
  \text{Audio} \xrightarrow{\text{VAD}} \xrightarrow{\text{MFCC}}
  \xrightarrow{\text{HMM}} \text{command}
\]

The VAD (voice activity detection) uses an adaptive energy threshold; each
segment is processed with LPC order~12, converted to Line Spectral
Frequencies (LSF), and scored against each per-word HMM. The word with the
highest log-likelihood wins.

The node publishes recognized tokens to \texttt{/voice\_command}; the
\texttt{stop} and \texttt{pause} tokens are consumed directly by the YASMIN
blackboard.

% ── IMAGE 3 ─────────────────────────────────────────────────────────────────
% Place: img/voice_pipeline.png
% Content: Confusion matrix of the HMM recognizer evaluated on the
%          11-word vocabulary.  Each row = true label, each column =
%          predicted label.  Diagonal should show per-word accuracy.
%          Generate with the test script in voice_control/scripts/.
\begin{figure}[htbp]
  \centering
  \includegraphics[width=\columnwidth]{img/voice_pipeline.jpeg}
  \caption{Confusion matrix of the LPC + HMM recognizer on the 11-word
           vocabulary. Diagonal cells show per-word recognition accuracy.}
  \label{fig:voice_confusion}
\end{figure}

\subsection{MR. Reference Detection --- ArUco }

The ArUco detector (\texttt{aruco\_detector.py}) was finalized in Sprint~3
and remained unchanged in Sprint~4. It is compatible with OpenCV~4.6 and
4.7+ through a version-detection shim applied at import time. Each detected
marker's full 3D pose (rvec + tvec → quaternion) is published to
\texttt{/aruco\_poses} (\texttt{PoseArray}).

\textbf{Pending:} the bridge that converts \texttt{/aruco\_poses} to
\texttt{PoseWithCovarianceStamped} on \texttt{/pose\_measurement} (the EKF
input) was not completed in this sprint.

\subsection{MR. Localization --- EKF + ICP }

The EKF node from Sprint~2 was kept intact. The ICP improvements in
Sprint~4 (better keyframe policy, outlier rejection) translate directly into
more accurate pose corrections published to \texttt{/pose\_measurement},
which the EKF fuses in its update step.

\textbf{Status:} EKF + ICP localization is functional on the real robot.
ArUco absolute corrections remain pending (bridge node not implemented).

\subsection{MR. Navigation}

The navigation stack received one significant addition in Sprint~4:

\paragraph{Local costmap} \texttt{local\_costmap.py} implements
a world-anchored, rolling LiDAR costmap with temporal decay:

\begin{itemize}[noitemsep, leftmargin=*]
  \item A square patch of the world, centred on the robot, stores float
        scores in $[0, 1]$ per cell.
  \item Scan hits raise the cell score to 1.0.
  \item Each control cycle applies exponential decay:
        $\text{score} \leftarrow \text{score} \cdot e^{-\Delta t/\tau}$,
        so transient detections (a person walking past) fade within $\sim$1~s
        while persistent obstacles stay occupied.
  \item The patch re-centres on demand; cells that leave the patch are
        dropped and new cells start at~0.
\end{itemize}

The local costmap feeds the Bug1 avoidance logic and the stuck detector in
\texttt{nav\_node.py}, replacing direct polling of the raw \texttt{/scan}
topic with a smoothed, decayed occupancy signal.

% ── IMAGE 4 ─────────────────────────────────────────────────────────────────
% Place: img/local_costmap.png
% Content: RViz2 screenshot showing the local costmap (colored grid patch)
%          overlaid on the global occupancy grid while the robot navigates.
%          Show at least one obstacle with its decay trail.
\begin{figure}[htbp]
  \centering
  \includegraphics[width=\columnwidth]{img/local_costmap.png}
  \caption{Local costmap overlay during navigation in Gazebo. The rolling
           patch is centred on the robot; decaying scores (lighter cells)
           mark the trail of a transient obstacle.}
  \label{fig:costmap}
\end{figure}

The \texttt{nav\_node.py} was also significantly refactored:
replanning on stuck detection, goal tolerance parameterisation, and
integration of the local costmap signal for reactive avoidance.
\texttt{navigation\_real.launch.py} was added as the real-robot
navigation launch file.

\subsection{MR. Target Detection and Alignment}

This subsystem received the most development effort of the sprint. Three
independent QR-based visual alignment strategies were implemented in the
\texttt{perception} package:

\paragraph{Strategy 1: Scripted state machine (\texttt{qr\_alignment.py})}

States: \texttt{IDLE → SEARCH → CENTER\_QR → ALIGN → DOCK → DONE}.
The robot first centres the QR marker in the camera frame (correcting heading
only), then aligns perpendicular to the QR without advancing (maintaining
distance $\geq D_{\text{dock}} + \text{margin}$ by forcing $v=0$ when too
close), and finally docks straight in. An exponential moving average (EMA)
filter on the PnP geometry reduces jitter from the 30~mm QR at 320$\times$240
resolution.

\paragraph{Strategy 2: Cubic Bézier + pure pursuit
(\texttt{qr\_path\_alignment.py})}

At every 10~Hz tick with a fresh detection:
\begin{enumerate}[noitemsep, leftmargin=*]
  \item Compute the goal pose in the robot frame: a point $D_{\text{dock}}$
        in front of the QR along its normal, facing the QR.
  \item Plan a cubic Bézier curve from $(0,0,0)$ to the goal with tangents
        that respect the initial and final headings.
  \item Pure pursuit: find the look-ahead point on the curve and command
        $v$ and $\omega$ to track it.
\end{enumerate}
Continuous replanning makes the path adapt to detection jitter without
discrete state transitions.

\paragraph{Strategy 3: Parabola + wheel odometry
(\texttt{qr\_quad\_alignment.py})}

\begin{enumerate}[noitemsep, leftmargin=*]
  \item \textbf{PLAN:} at detection, snapshot the QR pose. Compute the goal
        and plan a parabola $y = ax^2$ in the robot frame. Reset wheel
        odometry.
  \item \textbf{FOLLOW:} pure pursuit on the parabola using
        \textit{only wheel odometry} (from \texttt{odometry\_tracker.py}).
         No QR dependence — robust to occlusion during travel.
  \item \textbf{DOCK:} final straight approach with QR feedback if visible,
        open-loop otherwise.
\end{enumerate}

\paragraph{Camera calibration (\texttt{camera\_calibration.py})}

A ROS~2 node for interactive checkerboard calibration was implemented.
It subscribes to \texttt{/cam\_img}, accumulates corner detections with
SPACE, and computes the camera matrix $\mathbf{K}$ and distortion
coefficients $\mathbf{d}$ via OpenCV Zhang's method. Results are written to
\texttt{config/camera\_params.yaml} in the format expected by the QR pose
detector.

\paragraph{QR pose detector (\texttt{qr\_pose\_detector.py})}

Loads the calibrated intrinsics, decodes QR codes with \texttt{pyzbar},
estimates the 3D pose (solvePnP with the four corner correspondences), and
publishes to \texttt{/aruco\_poses} (\texttt{PoseArray}) — the same topic
as the ArUco detector, so the state machine consumes a single alignment
topic regardless of marker type.

% ── IMAGE 5 ─────────────────────────────────────────────────────────────────
% Place: img/qr_alignment.png
% Content: Four-panel figure:
%   (a) camera frame showing the detected QR with its axes overlay (PnP result)
%   (b) top-down diagram of the planned Bézier / parabola path from robot
%       to docking pose
%   (c) terminal screenshot or plot showing alignment error converging to 0
%   (d) photo of the robot docked in front of the QR-marked pallet position
\begin{figure}[htbp]
  \centering
  \includegraphics[width=\columnwidth]{img/qr_alignment.jpeg}
  \caption{QR-based visual alignment. (a)~Detected QR with 3-D pose axes
           overlaid (PnP). (b)~Planned Bézier path from current robot pose
           to docking pose. (c)~Bearing and lateral error converging to zero
           during a FOLLOW→DOCK sequence.}
  \label{fig:qr_align}
\end{figure}

\subsection{MR. Web Interface with Telemetry }

The Flask + SocketIO dashboard was completed in Sprint~3 and extended in
Sprint~4 with two new capabilities:

\paragraph{Live map rendering} \texttt{ros\_bridge.py} was extended to subscribe to \texttt{/map} (\texttt{OccupancyGrid}), encode each
update as a PNG, and push it to the frontend via a SocketIO
\texttt{map\_update} event. The frontend renders it on an HTML5 canvas,
overlaying the robot pose marker. This replaced the earlier static-map
approach and gives operators real-time SLAM feedback without RViz~2.

\paragraph{Voice-command dispatch} the frontend received a voice-command
panel: a microphone button records audio directly in the browser, sends it
to the Flask \texttt{/api/voice} endpoint as WebM, which is decoded with
\texttt{pydub} + \texttt{ffmpeg}, resampled to 16~kHz, and forwarded to
the HMM recognizer. The recognized token is published to
\texttt{/voice\_command}.

\paragraph{Final architecture}
\begin{itemize}[noitemsep, leftmargin=*]
  \item \textbf{Backend:} \texttt{dashboard\_node.py} subscribes to all
        telemetry topics and forwards them via SocketIO.
        \texttt{ros\_bridge.py} publishes to \texttt{/cmd\_vel} and
        \texttt{/mission} from Flask routes.
  \item \textbf{Frontend:} canvas map, MJPEG camera feed, telemetry panel
        (pose, YASMIN state, active mission), teleoperation controls, and
        voice-command microphone.
  \item \textbf{Safety watchdog:} 350~ms timer publishes zero velocity if
        no teleoperation commands arrive (dead man's switch).
\end{itemize}

\begin{figure}[htbp]
  \centering
  \fbox{\parbox{0.9\columnwidth}{\centering\vspace{4pt}
    \textbf{All Sprint~4 demo videos} \\[4pt]
    Mapping, navigation, QR alignment, lifter servo, dashboard live feed,
    and voice command demos. \\[6pt]
    \href{https://drive.google.com/drive/folders/1CV6YNE6Rj2bZOIr3SCd7zk2VpTP_-eQ4?usp=sharing}{%
      \texttt{drive.google.com/drive/folders/1CV6YNE6Rj2bZOIr3SCd7zk2VpTP\_-eQ4}} \\[4pt]
  }}
  \caption{Google Drive folder with all Sprint~4 demonstration videos.}
  \label{fig:all_videos}
\end{figure}

The source code for all packages is at
\url{https://github.com/Hugo734/Puzzlebot-Challenge-}

% ════════════════════════════════════════════════════════════════════════════
\section{Final System Architecture}

Table~\ref{tab:packages} summarises the eight packages that make up the
final system after Sprint~4.

\begin{table}[h]
\caption{ROS~2 packages in the final system}
\label{tab:packages}
\renewcommand{\arraystretch}{1.2}
\centering
\begin{tabular}{>{\raggedright}p{2.2cm}
                >{\centering}p{1.1cm}
                >{\raggedright\arraybackslash}p{4.0cm}}
\toprule
\textbf{Package} & \textbf{Language} & \textbf{Purpose} \\
\midrule
\texttt{localization}    & C++    & EKF + ICP + velocity bridge \\
\texttt{slam}            & Python & OccupancyGrid + saved map \\
\texttt{navigation}      & Python & A* + Bug1 + PID + local costmap \\
\texttt{mission\_control} & Python & YASMIN state machine \\
\texttt{perception}      & Python & ArUco + QR alignment + pallet \\
\texttt{lifting}         & Python & HAL GPIO → FPGA (3 bits) \\
\texttt{voice\_control}  & Python & LPC + HMM + parser \\
\texttt{dashboard}       & Python & Flask + SocketIO + voice API \\
\bottomrule
\end{tabular}
\end{table}

% ════════════════════════════════════════════════════════════════════════════
\section{Challenges Overcome}

\begin{itemize}[leftmargin=*]
  \item \textbf{Real-world SLAM smearing:} three independent root causes
        (DDS multicast collision, clock offset rejecting TF frames, QoS
        mismatch silently dropping scans) all produced the same symptom of a
        static map. Each was diagnosed separately and fixed in the real-world
        launch file.
  \item \textbf{QR PnP jitter at close range:} a 30~mm QR at 320$\times$240
        occupies only $\sim$30 pixels; the solvePnP solution jumps by several
        centimetres between frames. An EMA filter on the geometric parameters
        (bearing, distance, face angle) reduced the noise sufficiently for
        stable PID control.
  \item \textbf{Alignment during occlusion:} the Bézier strategy failed when
        the robot's body occluded the QR at close range. The parabola +
        odometry strategy (\texttt{qr\_quad\_alignment.py}) was developed
        specifically to address this: odometry drives the FOLLOW phase without
        requiring the QR to be visible.
  \item \textbf{HMM model path after \texttt{pip install}:} the Sprint~3
        recognizer searched a hardcoded relative path that only exists in the
        source tree, not after installation. Switching to
        \texttt{importlib.resources} made the node portable.
  \item \textbf{YASMIN + ROS~2 thread safety:} the state machine runs in a
        background thread; ROS~2 callbacks update the blackboard from the
        executor thread. Blackboard writes are atomic for all types used
        (str, float, dict) in CPython, which is sufficient for this use case.
  \item \textbf{Teleoperation watchdog:} a browser disconnect left the robot
        moving indefinitely. The 350~ms backend timer that publishes zero
        velocity acts as a dead man's switch.
\end{itemize}

% ════════════════════════════════════════════════════════════════════════════
\section{Conclusion and Next Steps}

Sprint~4 advanced the system from a simulation-validated prototype to a
hardware-ready stack. Real-world mapping now produces clean occupancy grids
on the Jetson Nano; the QR alignment pipeline provides three strategies at
increasing robustness; the dashboard gives operators live map feedback and
voice dispatch; and the Forklift FSM is modular and independently versioned.

\textbf{Remaining work for the final demonstration:}
\begin{itemize}[noitemsep, leftmargin=*]
  \item Record and label real waypoints on the competition track.
  \item Select and tune the best QR alignment strategy on hardware.
  \item Train the HMM codebooks with recordings in the noisy lab environment.
  \item Implement the ArUco--EKF bridge node for absolute pose corrections.
  \item Complete and validate the CNN trailer detector.
  \item Run the full pick-and-place sequence on the physical robot.
\end{itemize}

% ════════════════════════════════════════════════════════════════════════════
\section*{Sprint 4 Backlog}

\begin{table}[h]
\caption{Sprint 4 Task Status}
\label{tab:backlog}
\renewcommand{\arraystretch}{1.2}
\centering
\begin{tabular}{>{\raggedright}p{3.5cm}
                >{\centering}p{1.3cm}
                >{\centering}p{1.1cm}
                >{\centering\arraybackslash}p{1.0cm}}
\toprule
\textbf{Task} & \textbf{Assigned} & \textbf{SP} & \textbf{Status} \\
\midrule
Full system integration test               & Jordan    & 5 & In progress \\
Digital twin                               & Rosendo   & 3 & In progress \\
Vision precision validation                & Juan José & 3 & In progress \\
Gripper testing with real boxes            & Jordan    & 3 & In progress \\
Technical API documentation                & Rosendo   & 3 & Done \\
Global + local navigation integration      & Jordan    & 5 & Done \\
Detection--alignment pipeline (QR)         & Jordan    & 3 & Done \\
ISO 10218 \& 13482 documentation           & Hugo      & 2 & Done \\
YASMIN state machine                       & Jordan    & 5 & Done \\
Flask + SocketIO dashboard                 & Rosendo   & 5 & Done \\
HMM voice module + ROS~2 node             & Rosendo   & 5 & Done \\
A* + Bug1 + PID navigation                 & Jordan    & 8 & Done \\
Gazebo + RViz2 simulation integrated       & Jordan    & 3 & Done \\
Real-world mapping stack                   & Jordan    & 5 & Done \\
Local costmap                              & Jordan    & 3 & Done \\
Camera calibration                         & Jordan    & 2 & Done \\
Forklift FSM submodule                     & Hugo      & 3 & Done \\
Lifter URDF + mesh finalisation            & Juan José & 3 & Done \\
Lifter redesign + servo + FPGA demo        & Hugo      & 5 & Done \\
Integrate ArUco with EKF                   & Juan José & 3 & In progress \\
PID alignment with pallet                  & Victor    & 3 & In progress \\
CNN trailer detection                      & Juan José & 5 & In progress \\
\bottomrule
\end{tabular}
\end{table}

% ════════════════════════════════════════════════════════════════════════════
\section*{Individual Reflections}

\subsection*{Juan José Jáuregui Barba --- A00836722}

My contribution in Sprint~4 focused on the simulation model and
hardware-ready URDF. I reorganised all lifter mesh files into a
\texttt{lifter/} subfolder, rewrote the lifter URDF xacro to correctly
reference the new mesh paths (carriage, inner mast, outer mast, camera
holder), and updated \texttt{puzzlebot\_with\_lifter.urdf.xacro} with
proper joint limits and inertial properties. I also updated the warehouse
Gazebo world to add the rack model (\texttt{rack.dae}), which required
authoring the DAE file from scratch and tuning the collision geometry in
\texttt{model.sdf}. On the launch side, I removed obsolete files and
consolidated all simulation bringup into the updated \texttt{gazebo.launch.py}.
On Sprint~4 remaining tasks, the ArUco--EKF bridge node is my responsibility;
I understand what it needs (convert \texttt{PoseArray} to
\texttt{PoseWithCovarianceStamped} with a tuned covariance) and plan to
complete it before the final demonstration.

\subsection*{Victor Alejandro Meneses Garza --- A01384002}

I focused on the mechanical and electronic integration of the physical robot.
On the mechanical side, I secured the LiDAR at the correct height and
confirmed that the TF offset in the URDF (\texttt{lidar\_joint} origin)
matches the measured position on the chassis. On electronics, I reviewed the
H-bridge wiring, power distribution between the Jetson, motors, and FPGA,
and verified that the GPIO voltage levels (Jetson 3.3~V) are compatible with
the FPGA inputs. I also began the alignment PID controller for pallet
detection: pallet centroid detection via HSV segmentation is functional and
the pixel error is published to \texttt{/alignment\_error}, but the PID gains
$K_p$, $K_i$, $K_d$ have not yet been tuned on the real robot. My priority
for the final is to finish the gain tuning and validate the alignment sequence
end-to-end on hardware.

\subsection*{Jordan Palafox Salinas --- A00835705}

Sprint~4 was the most intensive sprint for me. I debugged and fixed the
entire real-world mapping stack (DDS, clock offset, QoS mismatch, TF
pipeline), which took two full days of hardware testing before the Jetson
produced a clean map. I then implemented the local costmap module
(\texttt{local\_costmap.py}) to give navigation a smoothed obstacle signal
with temporal decay. The largest effort was the QR alignment system: I
implemented three complete strategies---scripted state machine, Bézier +
pure pursuit, and parabola + wheel odometry---and the associated camera
calibration and QR pose detection nodes. The parabola + odometry approach
emerged from the failure of the vision-only strategies when the robot's body
occluded the QR at close range; switching to wheel odometry for the follow
phase was the key insight. I also consolidated \texttt{sim.launch.py} and
\texttt{real.launch.py} and added the map rendering to the dashboard.

\subsection*{Hugo Daniel Castillo Ovando --- A00836025}

My work in Sprint~4 covered three areas: the Forklift FSM, the new lifter
mechanical design, and ISO documentation. I packaged the Forklift
finite-state machine as an independent git submodule (\texttt{src/forklift\_fsm}),
which allows the FPGA team to version their firmware independently while the
ROS~2 stack consumes it through the \texttt{lifting\_node} interface.
On hardware, I led the redesign of the lifter mechanism: the new assembly
uses a servo motor controlled by the FPGA 3-bit output to reach the eight
discrete fork heights. The design was modelled in CAD, fabricated, and
validated end-to-end — a video of the Jetson publishing \texttt{/lifter\_level}
commands and the servo responding through all eight levels is linked in
Figure~\ref{fig:lifter_demo}. On ISO documentation, I mapped ISO~10218
and ISO~13482 requirements to the specific software mechanisms in the stack
(speed cap, protective stop, dead man's switch).

\subsection*{Rosendo De Los Rios Moreno --- A01198515}

In Sprint~4 I completed the voice recognition integration and the two
dashboard extensions. On voice, I fixed the HMM model path resolution
bug that broke the node after \texttt{pip install}, trained and committed
eleven per-word \texttt{.npz} models, and wired the \texttt{/api/voice}
Flask endpoint to decode WebM audio from the browser (using \texttt{pydub}
+ \texttt{ffmpeg}), resample it to 16~kHz, and pass it to the HMM
recognizer. On the dashboard, I added the live map rendering pipeline:
the ROS~2 \texttt{ros\_bridge.py} node subscribes to \texttt{/map},
encodes each \texttt{OccupancyGrid} update as a PNG, and pushes it via
SocketIO; the frontend renders it on a canvas. The result is that
operators have real-time SLAM feedback in the browser without needing
RViz~2. Remaining work: retraining the HMM codebooks with recordings
from the noisy lab environment, and adding a dashboard alert when the
YASMIN machine enters \texttt{STUCK}.

\end{document}
